\section{GPU実装}
\label{sec:gpu}

\subsection{想定ハードウェアとデータレイアウト}

本プロジェクトは単一のコンシューマGPU（NVIDIA GeForce RTX 3060 Ti，VRAM 8\,GB GDDR6）を実行環境として想定する。LBMはメモリ帯域律速のアルゴリズムであるため，実効性能はデータセンター向けGPUよりも一桁近く低く，8\,GBのVRAMが格子点数の上限を決める最大の制約となる。

実装言語はJuliaとし，GPUカーネルはCUDA.jlを用いて直接記述する。データレイアウトはSoA（Structure of Arrays）とし，分布関数 \(f_i\) を速度方向ごとに連続配列として保持することでメモリコアレッシングを最大化する。

\subsection{AA-pattern（in-place streaming）}

8\,GB VRAMという制約を踏まえ，AA-pattern（in-place streaming）を採用する（\texttt{src/core/aa\_pattern.jl}）。通常のping-pong方式（2バッファ）に対し，分布関数のメモリフットプリントを約半減できる：D3Q27・FP32では\(27\times4=108\,\text{byte}\)/ノードとなり，本番解像度（\(320^3\)）で約3.6\,GB。ping-pong方式なら約7.1\,GBとなり8\,GBのカードには収まらないため，この選択は必須である。

実装ではD3Q27を「立方体スロット順」で保持する：スロット \(s\) と速度 \(\vec c\) の対応を
\begin{align}
	c = \left((s-1)\bmod3,\ \left\lfloor\frac{s-1}{3}\right\rfloor\bmod3,\ \left\lfloor\frac{s-1}{9}\right\rfloor\right) - (1,1,1)
	\label{eq:cube-slot}
\end{align}
に固定することで，速度・重み・反対方向がすべて整数演算で求まり（反対スロットは\(28-s\)），カーネル内でのテーブル参照が不要になる。またAAパターンでは2つのノードが同じスロットに触れないため，スレッド間の競合が原理的に発生しない。

AAパターンで8ステップ進めた結果は，既存の2格子方式（解析解で検証済み）と相対誤差\(10^{-12}\)で一致することをBGK・中心モーメント両方で確認している。

\subsection{解像度とVRAMの見積り}
\label{sec:resolution-budget}

AA-pattern・D3Q27・FP32を前提に，1格子点あたりの必要メモリを見積もる。分布関数27方向で\(27\times4=108\,\text{byte}\)，固体フラグ等の付随データを含めて概算130--140\,byte程度になる。RTX 3060 Ti（8\,GB）のうちOS等のオーバーヘッドを見込み計算に使えるVRAMを約7\,GBと仮定すると，格子点数の上限は概ね\(5\times10^7\)点程度となる。

\begin{table}[H]
	\centering
	\caption{一様格子（立方体領域，N点/D）における解像度とコストの関係。}
	\label{tab:resolution-budget}
	\begin{tabular}{cccccccc}
		\Hline{1.5pt}
		N/D & L/D & 一辺 & VRAM & \(\Delta x\) (mm) & 縫い目/\(\Delta x\) & 境界層/\(\Delta x\) & 時間 \\
		\hline
		40  & 8 & 320 & 3.42\,G & 1.870 & \textbf{0.42} & 0.40 & 0.7\,h \\
		50  & 6 & 300 & 2.82\,G & 1.496 & 0.53 & 0.51 & 0.8\,h \\
		60  & 6 & 360 & 4.87\,G & 1.247 & 0.63 & 0.61 & 1.6\,h \\
		80  & 4 & 320 & 3.42\,G & 0.935 & 0.84 & 0.81 & 1.5\,h \\
		120 & 3 & 360 & 4.87\,G & 0.623 & 1.27 & 1.21 & 3.2\,h \\
		\hline
	\end{tabular}
\end{table}

\textbf{縫い目高さと境界層厚は同じ数字である。} 代表条件（\(U=39\,\text{m/s}\)）の赤道位置での層流境界層厚は\(5\sqrt{\nu x/U}\approx0.76\,\text{mm}\)であり，縫い目高さ0.79\,mmと5\%以内で一致する。粗さ要素が境界層と同程度の高さを持つことは，それが遷移を誘発するトリップとして働く条件そのものであり（\secref{aerodynamics}），同時に「縫い目を解像すること」と「境界層を解像すること」が2つの独立な要求ではなく1つであることを意味する。\tabref{resolution-budget}の縫い目/\(\Delta x\)列と境界層/\(\Delta x\)列が常に一緒に動くのはこのためである。

一様格子で領域を8Dに保つ限り，VRAM上限内で到達できる解像度は縫い目高さの0.42セル程度にとどまる。1セル以上に載せるには領域を3D程度まで狭めるしかないが，そこでは球が箱の1/6を占め，これはブロッケージ（\secref{validation}のHasimoto解比較で定量化した周期像の影響）そのものが無視できない大きさになってしまう。この一様格子の限界が，局所グリッドリファインメント（\secref{grid-refinement}）を最適化ではなく必須要件として位置づける根拠である。1レベル・2倍細分化の導入により，同じ8D領域・同程度のVRAMで縫い目解像度をおよそ2倍にできることを確認している。

\subsection{壁面境界のGPU実装}

融合カーネル（衝突・体積力・ストリーミングを1カーネルに融合したもの）では，バウンスバックを読み出し側ではなく\emph{書き出し側}に置く。書き出しの時点でスレッドは自ノードの衝突後分布を保持しており，それがバウンスバックに必要なものそのものだからである。跳ね返った値は，次のステップで自分自身が読みに来るスロットへ直接書き込む：偶数ステップで固体へ向かう方向 \(s\) は隣接固体ノードのスロットへ，奇数ステップでは自ノードの反対スロット \(28-s\) へ書く。これにより読み出し側には壁面の処理が一切不要になり，スレッド間の競合が生じないというAAパターンの性質もそのまま保たれる。跳ね返り値が同じステップ・同じノードの衝突後分布と密度から作られるため，2格子方式の参照実装と項ごとに一致し，近似ではなく厳密な比較ができる（分布関数で相対誤差\(10^{-12}\)，力・トルクで\(10^{-10}\)の一致を確認済み）。

境界ノードは自身の通し番号を持つため，各ノードが専用の列に力・トルクを書き出し，後段で総和を取る。アトミック演算は不要であり，2つのスレッドが同じ列に触れることもない。

この方式の制約として，Bouzidiの \(\delta<1/2\) 分岐は「1つ外側の流体ノードの衝突後分布」を必要とするが，融合カーネルのスレッドはそれを持たない。該当リンクはハーフウェイ規則にフォールバックする（\(\delta\ge1/2\)は自ノードの分布のみで済むため補間規則を使う）。2格子方式側にも同じ規則を用意してあるため参照実装と厳密に一致し，この妥協のコストは推測ではなく測定できる。

\subsection{局所グリッドリファインメント}
\label{sec:grid-refinement}

\subsubsection{音響スケーリング}

細分化比 \(m\) の局所グリッドリファインメントでは，時間刻み \(\Delta t\) を空間刻み \(\Delta x\) と一緒に半分にする（音響スケーリング）。時間刻みを両レベルで共通にすれば界面の時間補間が不要になり実装は大幅に楽になるが，その場合物理音速 \(c_{phys}=\Delta x/(\Delta t\sqrt3)\) が粗細で異なることになり，音波が界面でインピーダンス不整合を起こして部分反射する。移動境界は音響ノイズを出し続けるため，これを細格子内に閉じ込める必要がある。

\begin{table}[H]
	\centering
	\caption{細分化比 \(m\) に対する物理量のスケーリング。添字 c/f は粗/細格子を表す。}
	\label{tab:acoustic-scaling}
	\begin{tabular}{lll}
		\Hline{1.5pt}
		量 & 関係 & 理由 \\
		\hline
		緩和時間 & \(\tau_f-1/2=m(\tau_c-1/2)\) & \(\nu_{lat}=\nu\,\Delta t/\Delta x^2\) が \(m\) 倍 \\
		体積加速度 & \(a_f=a_c/m\) & \(a\propto\Delta t^2/\Delta x\) \\
		ひずみ速度 & \(S_f=S_c/m\) & \(S_{lat}=S_{phys}\Delta t\) \\
		格子速度・マッハ数 & 不変 & 音響スケーリングの定義 \\
		\hline
	\end{tabular}
\end{table}

非平衡部の再スケールは特に間違えやすい。分布関数はそのままでは移せない：平衡分布 \(f^{eq}\) は \(\rho,u\)のみに依存し，音響スケーリングはこれらを変えないためそのまま渡るが，非平衡部 \(f^{neq}\) はひずみ速度を担うため渡らない。Chapman-Enskog展開より \(f^{neq}\propto\tau_{lat}S_{lat}\) なので
\begin{align}
	\alpha \equiv \frac{f^{neq}_{fine}}{f^{neq}_{coarse}}
	= \frac{\tau_f}{m\,\tau_c}
	= \frac{\tau_f(\tau_c-1/2)}{\tau_c(\tau_f-1/2)}.
	\label{eq:alpha}
\end{align}
\(\alpha\)は\(\tau\)に依存する定数ではない——\(\tau\to1/2\)で\(1/m\)に，\(\tau\)が大きいと\(1\)に近づく。したがって\(\tau\)によらない式を用いると，投球レイノルズ数が強いる緩和時間域でちょうど\(m\)倍近く外れる一方，\(\tau\approx1\)のベンチマークでは数十\%の差にしか見えず離散化誤差と紛れる。

\subsubsection{配置と検証結果}

細格子は粗格子インデックスの範囲を覆う直方体パッチとし，一辺は\(2(\text{範囲})+1\)節点になるよう配置する。奇数インデックスが粗格子節点と一致し偶数が中間になるため，粗→細の転送に一般の補間は不要で，1・2・4・8点の平均で済む。1サイクルは粗2ステップ・細4ステップ（AAパターンの偶数レイアウトに両レベルが揃って戻る単位）である。

\begin{table}[H]
	\centering
	\caption{2格子リファインメントの検証結果。}
	\label{tab:refine-validation}
	\begin{tabular}{lc}
		\Hline{1.5pt}
		項目 & 結果 \\
		\hline
		一様流の転送（両方向） & 厳密（\(3\times10^{-18}\) / \(6\times10^{-17}\)） \\
		一様流の1サイクル（BGK） & 厳密（\(<10^{-14}\)） \\
		物理ひずみ速度の保存（粗↔細） & 相対\(10^{-11}\)以内 \\
		Taylor-Green渦の収束次数（音響スケーリング） & \textbf{1.98 / 2.42} \\
		\hline
	\end{tabular}
\end{table}

パッチ形状については，物体に沿った形（球殻など）ではなく軸並行の直方体を採用している。物体適合形状が有利に見えるのは節点数だけであり（同じ幅の球は立方体の52\%の体積），この節約は密な直方体配列から間接参照（索引付きリスト）へ格納形式を変えたときにのみ実体化する。実測した索引アクセスの帯域は直接コピーの0.11倍（局所性が最悪の場合）であり，実行全体で見込める節約の上限は11.7\%程度に対し，悪化の下限が4.8倍という非対称な結果から，直方体パッチを採用している。加えて，物体に密着した界面は剥離せん断層——本プロジェクトが最も精度を必要とする領域——の内部に置かれてしまうため，精度の観点からも直方体形状（界面を穏やかな流れの中に置く）が優れている。

パッチのアスペクト比については，ボールを中心とする立方体は前方の助走と後方の後流に同じ長さを割り当てるが，前方の流れは半径の半分も行けば滑らかである一方，後方の剥離とせん断層はどのレイノルズ数でもそれより長く伸びる。そこで各パッチを下流側にのみ延長するオプション（\texttt{--refine-wake}）を設けている。

\subsection{性能実測}

RTX 3060 Ti上での実測MLUPS（Mega Lattice Update Per Second）を示す。

\begin{table}[H]
	\centering
	\caption{精度・衝突演算子別の実測性能（周期境界のみ）。}
	\label{tab:mlups}
	\begin{tabular}{llrr}
		\Hline{1.5pt}
		精度 & 衝突演算子 & 実測MLUPS & 実効メモリ帯域 \\
		\hline
		FP32 & 中心モーメント & 2{,}544 & 549\,GB/s \\
		FP32 & BGK & 2{,}563 & 554\,GB/s \\
		FP64 & BGK & 246 & 106\,GB/s \\
		FP64 & 中心モーメント & 215 & 93\,GB/s \\
		\hline
	\end{tabular}
\end{table}

FP32ではカーネルはメモリ帯域に完全に張り付いている——BGKと中心モーメントが同一の性能を示すこと自体が帯域律速の証拠であり（演算量の差が性能に現れない），実測帯域（549--554\,GB/s）は同カードでの単純コピーベンチの実測帯域（513--532\,GB/s）とほぼ一致する。したがって，これ以上の高速化にはメモリトラフィック自体を減らす以外に手段がない。

一方，FP64は演算律速である（このカードのFP64は1/64レート）。BGKの方が中心モーメントより速いという順序が，演算律速であることを裏付けている。速度差が12倍あり，かつTaylor-Green渦での粘性再現はFP32でも誤差\(-0.06\%\)と精度上の問題がないため，本番計算はFP32とする。

壁面境界を含むカーネルの実測では，周期境界のみの場合に対する速度低下は3\%未満に留まる。境界に触れる節点が全体の1\%程度しかないうえ，バウンスバックを書き出し側に置いたことでギャザー側に壁の分岐が一切入らず，力・トルクのリダクションもアトミック操作を使わずに済むためである。

以上より，本番構成（8D領域・40点/D，\(320^3\)格子）での1条件あたりの計算時間は
\begin{align}
	\frac{3.28\times10^7\ \text{点}\times3.7\times10^5\ \text{ステップ}}{2.56\times10^9\ \text{更新/秒}} \approx 4.8\times10^3\ \text{秒} \approx \textbf{1.3時間/条件}
	\label{eq:runtime-estimate}
\end{align}
と見積もられる。これは計画段階の概算（15--20時間/条件）の約1/13であり，この余裕は解像度およびグリッドリファインメントに充てている。
